
%%%%%%%%%%%%%%%%%%%%%%%%%%%%%%%%%%%%%%%%%%%%%%%%%%%%%%%%%%%%%%%%%%%%%%%%%%%%%%%%%%%%%
%                                                                                   %
%     S.E. h�brido de clasificaci�n  y verificaci�n de originalidad de cl�sicos     %
%     ERL, Jes�s Bonilla                                                            %
%                                                                                   %
%     25-IV-2013                                                                    %
%%%%%%%%%%%%%%%%%%%%%%%%%%%%%%%%%%%%%%%%%%%%%%%%%%%%%%%%%%%%%%%%%%%%%%%%%%%%%%%%%%%%%

\documentclass[a4paper]{article} %
\usepackage{graphicx,amssymb} %

\usepackage{url}

\textwidth=15cm \hoffset=-1.2cm %
\textheight=25cm \voffset=-2cm %

\pagestyle{empty} %

\date{} %

\def\keywords#1{\begin{center}{\bf Keywords}\\{#1}\end{center}} %

\def\titulo#1{\title{#1}} %
\def\autores#1{\author{#1}} %

% Please, do not change any of the above lines

\begin{document}

% Type down your paper title
\titulo{A hybrid expert system for classic car\\
        recognition and originality evaluation}
% Authors
\autores{Eugenio Roanes-Lozano \\ %  Affiliation
       Instituto de Matem\'atica Interdisciplinar (IMI),\\
       Departamento de \'Algebra, Facultad de Educaci\'on, \\
       Universidad Complutense de Madrid, E-28040 Madrid (Spain) \\ \\
       Jes\'us Bonilla\\
       Motor Cl\'asico magazine (Editor in Chief)\\
       c/ Ancora 40, E-28045 Madrid (Spain)\\ \\
       \tt{eroanes@mat.ucm.es} % jbonilla@mpib.es} % Only one corresponding e-mail
       }%


\maketitle

\thispagestyle{empty}



% The abstract

\begin{abstract}
Apart from the condition of the car, one of the main problems when
buying a classic car (or when estimating its value) is recognizing
its originality.

Many times, non-original parts and accessories have been
installed, what is not normally very difficult to find out. But
sometimes items from other versions of the same model can be
found in a certain specimen, what is not that easy to recognize.

We have developed in the past Rule Based Expert Systems (RBES) and
AI tools for decision taking in different fields (medicine,
transportation engineering,...), both using algebraic inference
engines and logic programming.

The key idea of this work is to develop a computer package that
guides the user along a sequence of questions, in order to find
out the model and/or version of the car and to detect non-original
elements.

One important fact is that the available information can be
incomplete (for instance, all the available information about the
vehicle can be a set of photographs). Moreover, the user of the
system can be sometimes unable to answer all questions, even with
all the required information.

We have decided to use a hybrid approach, because:
\begin{itemize}
\item on one hand there are ``conclusive items'' that are easier to
      handle using rules in the style of classic RBES,
\item on the other hand, when the vehicle has a mixture of
      characteristics from different models/versions, it is easier
      to compare them with the predefined ones (stored in row
      matrix, sequence, list format,...) in order to identify
      the ``closest'' model/version.
\end{itemize}

A computer algebra system offers all the necessary tools for
implementing such an approach. The article is illustrated with a
small system (implemented in the computer algebra system
\textit{Maple}, that takes advantage of its \textit{Logic}
package), devoted to the Porsche 928.

\end{abstract}

\keywords{Expert Systems; Pattern Matching; Classic Cars;
          Computer Algebra Systems} % Write down at least 3 Keywords

% \section{Introduction}

\end{document}
